% Semantic tensor-network TikZ grammar.
%
% This file fixes the meaning of every glyph, wire, and region used in the
% blueprint.  Reusable MPS, MPO, MPDO, and PEPS constructions belong in
% macros/tn_library.tex; complete chapter-facing diagrams belong in
% macros/tn_print.tex.
%
% A solid line is always an index line.  Virtual and physical index lines have
% different widths.  Dashed curves are not contractions: they indicate a
% grouping, a basis change, or an identified periodic boundary.  Tensor labels
% do not change the glyph used for the tensor.
\usetikzlibrary{backgrounds,calc,fit}

\tikzset{
    tn picture/.style={
            baseline=-0.1cm,
            transform shape,
            font=\scriptsize
        },
    % Semantic node classes.
    tn tensor node/.style={
            draw,
            fill,
            shape=circle,
            inner sep=0.055cm
        },
    tn insertion node/.style={
            tn tensor node,
            red
        },
    % A junction is a contraction point of index lines, not a local tensor.
    tn junction node/.style={
            fill,
            shape=circle,
            inner sep=0.020cm
        },
    tn component node/.style={
            draw,
            fill=white,
            inner sep=1.4pt
        },
    tn factor node/.style={
            draw=black!65,
            fill=black!5,
            inner sep=1.4pt
        },
    tn map node/.style={
            draw,
            rounded corners=2pt,
            fill=blue!8,
            inner sep=2pt
        },
    tn state node/.style={
            draw=black!65,
            rounded corners=2pt,
            fill=black!4,
            inner sep=2pt
        },
    % Semantic index lines.
    tn virtual wire/.style={line width=0.4pt},
    tn physical wire/.style={line width=0.7pt},
    tn virtual trace/.style={tn virtual wire, rounded corners=3pt},
    tn physical trace/.style={tn physical wire, rounded corners=3pt},
    tn bent wire/.style={tn virtual wire, rounded corners=3pt},
    % Regions and non-contraction marks.
    tn factor region/.style={tn virtual wire, rounded corners=2pt},
    tn grouping region/.style={densely dashed, rounded corners=2pt},
    tn basis change mark/.style={densely dashed, line width=0.4pt},
    tn periodic mark/.style={tn virtual wire, densely dashed, ->},
    tn morphism arrow/.style={->, line width=0.60pt},
    tn comparison arrow/.style={<->, line width=0.60pt},
    tn selected path/.style={red, line width=0.70pt},
    tn secondary path/.style={blue, line width=0.70pt},
    tn region selected/.style={
            draw=red!75!black,
            fill=red!28,
            thick,
            rounded corners=1mm
        },
    tn region secondary/.style={
            draw=blue!75!black,
            fill=blue!24,
            thick,
            rounded corners=1mm
        },
    tn region complement/.style={
            draw=black!55,
            fill=black!8,
            densely dashed,
            thick,
            rounded corners=1mm
        },
    tn region collar/.style={
            draw=purple!70!black,
            fill=purple!18,
            thick,
            rounded corners=1mm
        },
    tn edge distinguished/.style={line width=0.6mm, green!65!black},
    tn vertex distinguished/.style={
            circle,
            draw=black!65,
            fill=yellow!75!white,
            thick
        },
    tn region label selected/.style={text=red!75!black},
    tn region label secondary/.style={text=blue!75!black},
    tn region label complement/.style={text=black!65},
    tn region label collar/.style={text=purple!70!black}
}

\makeatletter

\def\TN@physleglen{0.42}
\def\TN@physlabeloff{0.68}
\def\TN@doublelayersep{0.50}
\def\TN@chainsep{1.45}
\def\TN@chainright{4.20}
\def\TN@chainellipsisx{2.82}

% ---------- Semantic nodes ----------

\newcommand{\TN@node}[5][]{%
    \node[#2,#1] (#3) at #4 {#5};%
}

\newcommand{\TN@tensor}[3][]{%
    \TN@node[#1]{tn tensor node}{#2}{#3}{}%
}

% Anonymous tensor sites are used for lattice backgrounds whose sites are not
% referenced individually by the surrounding construction.
\newcommand{\TN@anonymoustensor}[2][]{%
    \node[tn tensor node,#1] at #2 {};%
}

\newcommand{\TN@component}[4][]{%
    \TN@node[#1]{tn component node}{#2}{#3}{$#4$}%
}

\newcommand{\TN@factor}[4][]{%
    \TN@node[#1]{tn factor node}{#2}{#3}{$#4$}%
}

\newcommand{\TN@map}[4][]{%
    \TN@node[#1]{tn map node}{#2}{#3}{$#4$}%
}

\newcommand{\TN@state}[4][]{%
    \TN@node[#1]{tn state node}{#2}{#3}{$#4$}%
}

\newcommand{\TN@insertion}[3][]{%
    \TN@node[#1]{tn insertion node}{#2}{#3}{}%
}

\newcommand{\TN@junction}[3][]{%
    \TN@node[#1]{tn junction node}{#2}{#3}{}%
}

\newcommand{\TN@anonymousjunction}[2][]{%
    \node[tn junction node,#1] at #2 {};%
}

\newcommand{\TN@scalar}[4][]{%
    \TN@tensor[#1]{#2}{#3}%
    \node[anchor=west] at ($(#2.east)+(0.10,0)$) {$#4$};%
}

% An external label does not alter the glyph of the object it names.
\newcommand{\TN@label}[4][]{%
    \node[#1] at ($(#2)+#3$) {$#4$};%
}

% ---------- Named ports ----------

% Every composable object may give a stable name to one of its boundary
% anchors.  Contractions can then be expressed solely in terms of ports.
\newcommand{\TN@port}[3]{%
    \coordinate (#1) at (#2.#3);%
}

% A port may also lie at a specified fraction of a boundary segment.  This is
% useful for a box carrying several indices on the same side.
\newcommand{\TN@betweenport}[4]{%
    \coordinate (#1) at ($(#2)!#4!(#3)$);%
}

\newcommand{\TN@westport}[3]{%
    \TN@betweenport{#1}{#2.north west}{#2.south west}{#3}%
}

\newcommand{\TN@eastport}[3]{%
    \TN@betweenport{#1}{#2.north east}{#2.south east}{#3}%
}

\newcommand{\TN@northport}[3]{%
    \TN@betweenport{#1}{#2.north west}{#2.north east}{#3}%
}

\newcommand{\TN@southport}[3]{%
    \TN@betweenport{#1}{#2.south west}{#2.south east}{#3}%
}

% ---------- Anchored index lines ----------

\newcommand{\TN@connectpoints}[4][]{%
    \draw[#2,#1] (#3) to (#4);%
}

\newcommand{\TN@vconnectpoints}[3][]{%
    \TN@connectpoints[#1]{tn virtual wire}{#2}{#3}%
}

\newcommand{\TN@pconnectpoints}[3][]{%
    \TN@connectpoints[#1]{tn physical wire}{#2}{#3}%
}

% Port-to-port composition is the preferred form when a construction is
% assembled from independently defined atoms.
\newcommand{\TN@vconnectports}[3][]{%
    \TN@vconnectpoints[#1]{#2}{#3}%
}

\newcommand{\TN@pconnectports}[3][]{%
    \TN@pconnectpoints[#1]{#2}{#3}%
}

\newcommand{\TN@vconnect}[5][]{%
    \TN@vconnectpoints[#1]{#2.#3}{#4.#5}%
}

\newcommand{\TN@pconnect}[5][]{%
    \TN@pconnectpoints[#1]{#2.#3}{#4.#5}%
}

% Curved port-to-port contractions retain the same compositional interface as
% straight contractions.  The two angles prescribe the tangent directions at
% the boundary ports.
\newcommand{\TN@vconnectcurve}[7][]{%
    \draw[tn virtual wire,#1] (#2.#3) to[out=#4,in=#5] (#6.#7);%
}

\newcommand{\TN@vconnectbendright}[6][]{%
    \draw[tn virtual wire,#1] (#2.#3) to[bend right=#4] (#5.#6);%
}

% Cubic connectors are useful when the route, rather than only its endpoint
% tangents, carries geometric information in the diagram.
\newcommand{\TN@vconnectbezier}[7][]{%
    \draw[tn virtual wire,#1] (#2.#3) .. controls #4 and #5 .. (#6.#7);% chktex 26
}

\newcommand{\TN@pconnectbezier}[7][]{%
    \draw[tn physical wire,#1] (#2.#3) .. controls #4 and #5 .. (#6.#7);% chktex 26
}

\newcommand{\TN@vtraceconnectbezier}[7][]{%
    \draw[tn virtual trace,#1] (#2.#3) .. controls #4 and #5 .. (#6.#7);% chktex 26
}

\newcommand{\TN@vbezieropen}[7][]{%
    \draw[tn virtual wire,#1] (#2.#3) .. controls #4 and #5 .. #6 -- ++(#7);% chktex 26
}

% A partial contraction terminates at the indicated fraction of the segment
% from the object centre to an external port.
\newcommand{\TN@vpartialconnect}[5][]{%
    \draw[tn virtual wire,#1] (#2.#3) -- ($(#2)!#5!(#4)$);%
}

% A trace may pass through an insertion lying below three consecutive atoms.
\newcommand{\TN@vtracevia}[6][]{%
    \draw[tn virtual trace,#1] (#2.west) -- ++(-#5,0) |- (#3.west);%
    \draw[tn virtual trace,#1] (#3.east) -| ($(#4.east)+(#6,0)$)
        -- (#4.east);%
}

\newcommand{\TN@vtraceconnectpoints}[3][]{%
    \TN@connectpoints[#1]{tn virtual trace}{#2}{#3}%
}

% A trace closure is itself composable when its two endpoints are named ports.
% The last two arguments give the horizontal and vertical clearances.
\newcommand{\TN@vtraceportsbelow}[5][]{%
    \draw[tn virtual trace,#1] (#2) -- ++(-#4,0) -- ++(0,-#5)
        -| ($(#3)+(#4,0)$) -- (#3);%
}

\newcommand{\TN@vtraceportsabove}[5][]{%
    \draw[tn virtual trace,#1] (#2) -- ++(-#4,0) -- ++(0,#5)
        -| ($(#3)+(#4,0)$) -- (#3);%
}

% General paths remain available for closures and lattice edges whose route is
% not a single segment.  Complete contractions should still use the anchored
% connect commands above whenever both endpoint objects are named.
\newcommand{\TN@vpath}[2][]{%
    \draw[tn virtual wire,#1] #2;%
}

\newcommand{\TN@ppath}[2][]{%
    \draw[tn physical wire,#1] #2;%
}

\newcommand{\TN@vtracepath}[2][]{%
    \draw[tn virtual trace,#1] #2;%
}

\newcommand{\TN@ptracepath}[2][]{%
    \draw[tn physical trace,#1] #2;%
}

\newcommand{\TN@arrow}[2][]{%
    \draw[tn morphism arrow,#1] #2;%
}

\newcommand{\TN@comparisonarrow}[2][]{%
    \draw[tn comparison arrow,#1] #2;%
}

\newcommand{\TN@selectedpath}[2][]{%
    \draw[tn selected path,#1] #2;%
}

\newcommand{\TN@secondarypath}[2][]{%
    \draw[tn secondary path,#1] #2;%
}

\newcommand{\TN@grouppath}[2][]{%
    \draw[tn grouping region,#1] #2;%
}

\newcommand{\TN@factorpath}[2][]{%
    \draw[tn factor region,#1] #2;%
}

\newcommand{\TN@openwire}[5][]{%
    \draw[#2,#1] (#3.#4) -- ++(#5);%
}

\newcommand{\TN@openport}[4][]{%
    \draw[#2,#1] (#3) -- ++(#4);%
}

\newcommand{\TN@vopenport}[3][]{%
    \TN@openport[#1]{tn virtual wire}{#2}{#3}%
}

\newcommand{\TN@popenport}[3][]{%
    \TN@openport[#1]{tn physical wire}{#2}{#3}%
}

\newcommand{\TN@vleg}[4][]{%
    \TN@openwire[#1]{tn virtual wire}{#2}{#3}{#4}%
}

\newcommand{\TN@pleg}[4][]{%
    \TN@openwire[#1]{tn physical wire}{#2}{#3}{#4}%
}

% Trace closures are index lines, not region boundaries.
\newcommand{\TN@vtracebelow}[3]{%
    \draw[tn virtual trace] (#1.west) -- ++(-0.42,0) -- ++(0,-#3)
        -| ($(#2.east)+(0.42,0)$) -- (#2.east);%
}

\newcommand{\TN@vtraceabove}[3]{%
    \draw[tn virtual trace] (#1.west) -- ++(-0.42,0) -- ++(0,#3)
        -| ($(#2.east)+(0.42,0)$) -- (#2.east);%
}

\newcommand{\TN@vtraceright}[3]{%
    \draw[tn virtual trace] (#1.east) -- ++(#3,0) |- (#2.east);%
}

\newcommand{\TN@ptraceright}[3]{%
    \draw[tn physical trace] (#1.north) -- ++(0,#3)
        -| ($(#1.center)!0.5!(#2.center)+(0.42,0)$)
        |- ($(#2.south)+(0,-#3)$) -- (#2.south);%
}

% ---------- Semantic regions ----------

\newcommand{\TN@region}[5][]{%
    \begin{scope}[on background layer]
        \node[#2, fit=#4, inner sep=4pt, #1] (#3) {#5};%
    \end{scope}%
}

\newcommand{\TN@groupregion}[4][]{%
    \TN@region[#1]{tn grouping region}{#2}{#3}{#4}%
}

\newcommand{\TN@factorregion}[4][]{%
    \TN@region[#1]{tn factor region}{#2}{#3}{#4}%
}

\newcommand{\TN@supportregion}[5][]{%
    \TN@region[#1]{tn region #2}{#3}{#4}{#5}%
}

% ---------- Index labels ----------

\newcommand{\TN@uplabel}[3]{%
    \node at ($(#1.north)+(0,#2)$) {$#3$};%
}

\newcommand{\TN@downlabel}[3]{%
    \node at ($(#1.south)+(0,-#2)$) {$#3$};%
}

\makeatother
